\paragrafos{Μέση τιμή και διάμεσος}{Γ.2.1, «Μέση τιμή, διάμεσος»}

\begin{anagke}
Η Ρήνη περιμένει κάθε πρωί το λεωφορείο για το σχολείο και μια εβδομάδα κρατά σημειώσεις: έχασε $5$, $6$, $4$, $5$, $7$, $5$ και $33$ λεπτά. Την Παρασκευή υπήρχε απεργία και το λεωφορείο καθυστέρησε πολύ. «Συνήθως περιμένω γύρω στα πέντε λεπτά», λέει η Ρήνη. Ο Ηλίας προσθέτει τους αριθμούς και τους μοιράζει στις μέρες: «Εγώ βρίσκω περίπου $9{,}3$ λεπτά». Ποιος έχει δίκιο;

Η κυρία Αλεξάνδρα ρωτά: «Και οι δύο έχετε δίκιο; Τι ακριβώς ψάχνει ο καθένας σας;». Επτά αριθμοί είναι δύσκολο να τους θυμάται κανείς. Χρειαζόμαστε \emph{έναν} αριθμό που να τους αντιπροσωπεύει και να μας λέει «πού βρίσκονται» τα δεδομένα. Αυτό όμως μπορεί να γίνει με δύο τρόπους.

Σε αυτή την παράγραφο μαθαίνουμε δύο μέτρα θέσης, τη μέση τιμή και τη διάμεσο, τα υπολογίζουμε από λίστα, από πίνακα συχνοτήτων και από ομαδοποιημένα δεδομένα, και καταλαβαίνουμε πότε προτιμούμε το ένα και πότε το άλλο.
\end{anagke}

\begin{lexi}{μέτρο θέσης, μέση τιμή, διάμεσος}{ονόματα και σημασίες}
Το \lexiword{μέτρο θέσης} δείχνει «πού βρίσκεται» το κέντρο των δεδομένων. Η \lexiword{μέση τιμή} λέγεται και \emph{μέσος όρος}: είναι ο αριθμός που θα είχαμε αν μοιράζαμε τα πάντα εξίσου. Η \lexiword{διάμεσος} είναι η τιμή που κάθεται «στη μέση» (δια-μέσος) όταν βάλουμε τα δεδομένα στη σειρά.
\end{lexi}

\begin{metafrasi}
Προσθέτω όλες τις τιμές και τις μοιράζω εξίσου. & μέση τιμή\\
Βάζω τις τιμές στη σειρά και παίρνω αυτή που είναι στη μέση. & διάμεσος\\
Μετρήσεις που κάπου τις έχω γράψει πολλές φορές. & πίνακας συχνοτήτων\\
Μετρήσεις που τις έχω χωρίσει σε κλάσεις (διαστήματα). & ομαδοποιημένα δεδομένα\\
\end{metafrasi}

\tmima{Η θεωρία}

\begin{kanonas}{μέση τιμή}
$$\text{μέση τιμή}=\frac{\text{άθροισμα των δεδομένων}}{\text{πλήθος των δεδομένων}}$$
\end{kanonas}

Η μέση τιμή είναι σαν σημείο ισορροπίας. Αν τοποθετήσεις τα δεδομένα $8, 9, 12, 15, 21$ πάνω σε μια αριθμογραμμή, σαν βαράκια σε μια ζυγαριά, το σημείο ισορροπίας βρίσκεται στο $13$: ο άξονας δεν γέρνει ούτε δεξιά ούτε αριστερά.

\begin{center}
\begin{tikzpicture}[x=0.5cm]
\draw[-{Stealth},thick] (6,0)--(23,0);
\foreach \x in {6,8,...,22}{\draw (\x,0.1)--(\x,-0.1) node[below,font=\scriptsize]{$\x$};}
\foreach \x in {8,9,12,15,21}{\fill[gblue] (\x,0.3) circle (3.5pt);}
\fill[gwarmD] (13,-0.55)--(12.65,-1.15)--(13.35,-1.15)--cycle;
\node[font=\small,gwarmD] at (13,-1.6) {μέση τιμή $=13$};
\end{tikzpicture}
\end{center}

\begin{idiotita}{τι ξέρουμε για τη μέση τιμή}
\begin{itemize}[leftmargin=1.2em,itemsep=0.3ex,topsep=2pt]
\item Είναι πάντα ανάμεσα στη μικρότερη και τη μεγαλύτερη τιμή.
\item Δεν είναι απαραίτητο να είναι μία από τις τιμές των δεδομένων (και μπορεί να είναι δεκαδικός).
\item Αν $\nu$ δεδομένα έχουν μέση τιμή $m$, τότε το άθροισμά τους είναι $\nu\cdot m$.
\item Την επηρεάζουν πολύ οι πολύ μεγάλες ή οι πολύ μικρές τιμές.
\end{itemize}
\end{idiotita}

\begin{kanonas}{διάμεσος}
Βάζω τα δεδομένα σε \textbf{αύξουσα σειρά}. Αν το πλήθος τους είναι \textbf{περιττό}, η διάμεσος είναι η \emph{μεσαία} τιμή. Αν το πλήθος τους είναι \textbf{άρτιο}, η διάμεσος είναι ο \emph{μέσος όρος των δύο μεσαίων} τιμών.
\end{kanonas}

Η διάμεσος χωρίζει τα δεδομένα στη μέση: τα μισά είναι μικρότερα ή ίσα με αυτήν και τα μισά μεγαλύτερα ή ίσα. Αν ένα δεδομένο γίνει πολύ μεγάλο, η διάμεσος σχεδόν δεν αλλάζει, ενώ η μέση τιμή «τραβιέται» προς αυτό.

\begin{kanonas}{μέση τιμή από πίνακα συχνοτήτων}
Όταν οι τιμές επαναλαμβάνονται, δεν προσθέτουμε μία μία. Εκτελούμε τρία βήματα:\\
1. Πολλαπλασιάζουμε κάθε τιμή με τη συχνότητά της.\\
2. Προσθέτουμε όλα τα γινόμενα.\\
3. Διαιρούμε με το σύνολο των παρατηρήσεων (το άθροισμα των συχνοτήτων).
\end{kanonas}

\begin{kanonas}{μέση τιμή από ομαδοποιημένα δεδομένα}
Όταν τα δεδομένα είναι χωρισμένα σε κλάσεις, δεν ξέρουμε την ακριβή τιμή κάθε παρατήρησης. Δεχόμαστε ότι όλες οι τιμές μιας κλάσης είναι ίσες με το \textbf{κέντρο της κλάσης} (το μέσο των άκρων της). Μετά κάνουμε τα ίδια τρία βήματα: κέντρο $\times$ συχνότητα, άθροισμα γινομένων, διαίρεση με το σύνολο. Το αποτέλεσμα είναι προσεγγιστικό.
\end{kanonas}

\begin{prosoxi}
\begin{itemize}[leftmargin=1.2em,itemsep=0.3ex,topsep=2pt]
\item Για τη διάμεσο \emph{πρώτα} βάζεις τις τιμές στη σειρά. Αν πάρεις την τιμή που κάθεται στη μέση της αταξινόμητης λίστας, κάνεις λάθος.
\item Στον πίνακα συχνοτήτων διαιρείς με το άθροισμα των \emph{συχνοτήτων}, όχι με το πλήθος των διαφορετικών τιμών.
\item Στα ομαδοποιημένα δεδομένα πολλαπλασιάζεις το \emph{κέντρο} της κλάσης με τη συχνότητα, όχι το αριστερό ή το δεξί άκρο.
\item Η διάμεσος δεν είναι ο μέσος όρος όλων των τιμών. Μόνο όταν το πλήθος είναι άρτιο παίρνουμε μέσο όρο, και μόνο των δύο μεσαίων.
\end{itemize}
\end{prosoxi}

\tmima{Λυμένα παραδείγματα}

\paradeigma{Η μέση τιμή}
\tixero Ο Νίκος σε πέντε αγώνες μπάσκετ πέτυχε $12$, $8$, $15$, $9$ και $21$ πόντους.\par
\tithelo Τη μέση τιμή των πόντων του.\par
\possindeo Προσθέτω τις πέντε τιμές και διαιρώ με το πλήθος τους, $5$.\par
\apantisi $\dfrac{12+8+15+9+21}{5}=\dfrac{65}{5}=13$ πόντοι ανά αγώνα.\par
\elegxos Το $13$ είναι ανάμεσα στο $8$ και στο $21$, όπως πρέπει. Και $5\cdot13=65$, που είναι το άθροισμα.

\paradeigma{Η διάμεσος, περιττό και άρτιο πλήθος}
\tixero Οι αναμονές της Ρήνης: $5, 6, 4, 5, 7, 5, 33$ λεπτά.\par
\tithelo Τη διάμεσο των επτά τιμών και τη διάμεσο αν χανόταν η καταγραφή της Δευτέρας ($5$), δηλαδή για τις $6, 4, 5, 7, 5, 33$.\par
\possindeo Βάζω τις τιμές στη σειρά. Επτά τιμές: μεσαία είναι η τέταρτη. Έξι τιμές: μέσος όρος της τρίτης και της τέταρτης.\par
\apantisi Επτά τιμές: $4, 5, 5, \underline{5}, 6, 7, 33$, άρα διάμεσος $5$ λεπτά. Έξι τιμές: $4, 5, \underline{5}, \underline{6}, 7, 33$, άρα διάμεσος $\dfrac{5+6}{2}=5{,}5$ λεπτά.\par
\elegxos Η μέση τιμή των επτά είναι $\dfrac{65}{7}\approx 9{,}3$ και την τραβά προς τα πάνω η τιμή $33$. Η διάμεσος $5$ δείχνει καλύτερα την «συνηθισμένη» αναμονή.

\paradeigma{Ποια τιμή λείπει;}
\tixero Ο Ηλίας έγραψε στα Μαθηματικά $15$, $17$, $14$, $18$ σε τέσσερα τεστ και θέλει μέσο όρο $16$ μετά το πέμπτο.\par
\tithelo Τι βαθμό χρειάζεται στο πέμπτο τεστ;\par
\possindeo Αν η μέση τιμή των $5$ τεστ είναι $16$, το άθροισμά τους είναι $5\cdot16=80$.\par
\apantisi Τα τέσσερα τεστ δίνουν $15+17+14+18=64$, άρα στο πέμπτο χρειάζεται $80-64=16$.\par
\elegxos $\dfrac{15+17+14+18+16}{5}=\dfrac{80}{5}=16$.

\paradeigma{Πίνακας συχνοτήτων}
\tixero Η Φανή ρώτησε $25$ συμμαθητές πόσα αδέλφια έχουν.\par
\tithelo Τη μέση τιμή.\par
\possindeo Προσθέτω μια στήλη «τιμή $\times$ συχνότητα».\par

\begin{center}
\renewcommand{\arraystretch}{1.25}
\begin{tabular}{lccccc}
\toprule
Αδέλφια & $0$ & $1$ & $2$ & $3$ & $4$\\
Συχνότητα & $3$ & $9$ & $8$ & $4$ & $1$\\
Τιμή $\times$ συχνότητα & $0$ & $9$ & $16$ & $12$ & $4$\\
\bottomrule
\end{tabular}
\end{center}

\apantisi Σύνολο παρατηρήσεων $3+9+8+4+1=25$. Άθροισμα γινομένων $0+9+16+12+4=41$. Μέση τιμή $\dfrac{41}{25}=1{,}64$ αδέλφια.\par
\elegxos Το $1{,}64$ είναι ανάμεσα στο $0$ και στο $4$. Δεν υπάρχει μαθητής με $1{,}64$ αδέλφια, αλλά ως μέσος όρος έχει νόημα.

\paradeigma{Ομαδοποιημένα δεδομένα}
\tixero Ο χρόνος που κάνουν $25$ μαθητές να φτάσουν στο σχολείο ομαδοποιήθηκε σε κλάσεις.\par
\tithelo Τον μέσο χρόνο διαδρομής.\par
\possindeo Βρίσκω το κέντρο κάθε κλάσης (π.χ. για $0$--$5$ είναι $\frac{0+5}{2}=2{,}5$) και συνεχίζω όπως στον πίνακα συχνοτήτων.\par

\begin{center}
\renewcommand{\arraystretch}{1.25}
\begin{tabular}{lccccc}
\toprule
Χρόνος (min) & $0$--$5$ & $5$--$10$ & $10$--$15$ & $15$--$20$ & $20$--$25$\\
Κέντρο κλάσης & $2{,}5$ & $7{,}5$ & $12{,}5$ & $17{,}5$ & $22{,}5$\\
Συχνότητα & $4$ & $9$ & $7$ & $3$ & $2$\\
Κέντρο $\times$ συχνότητα & $10$ & $67{,}5$ & $87{,}5$ & $52{,}5$ & $45$\\
\bottomrule
\end{tabular}
\end{center}

\apantisi Σύνολο $4+9+7+3+2=25$ μαθητές. Άθροισμα γινομένων $10+67{,}5+87{,}5+52{,}5+45=262{,}5$. Μέση τιμή $\dfrac{262{,}5}{25}=10{,}5$ λεπτά.\par
\elegxos Η μέση τιμή είναι ανάμεσα στο $0$ και στο $25$ και πλησιάζει τις κλάσεις με τις περισσότερες παρατηρήσεις ($5$--$10$ και $10$--$15$). Είναι προσέγγιση, γιατί δεν ξέρουμε τις ακριβείς τιμές.

\tmima{Ασκήσεις}

\zoni{Το κατάλαβα}{μέση τιμή, διάμεσος, πίνακας}

\ex Βρες τη μέση τιμή: α) $24$ παρατηρήσεις με άθροισμα $120$. β) $8$ παρατηρήσεις με άθροισμα $20$. γ) $40$ παρατηρήσεις με άθροισμα $10$.

\ex Βρες το άθροισμα: α) $30$ παρατηρήσεων με μέση τιμή $4$. β) $6$ παρατηρήσεων με μέση τιμή $2{,}5$. γ) $20$ παρατηρήσεων με μέση τιμή $0{,}1$.

\ex Βρες το πλήθος των παρατηρήσεων που έχουν: α) μέση τιμή $4$ και άθροισμα $48$. β) μέση τιμή $2{,}5$ και άθροισμα $30$. γ) μέση τιμή $0{,}5$ και άθροισμα $20$.

\ex Υπολόγισε τη μέση τιμή: α) $4, 6, 8, 10, 12$. β) $-2, 0, 3, 5$ (θερμοκρασίες σε $^\circ$C). γ) $7, 7, 7, 7$.

\ex Βρες τη διάμεσο: α) $9, 3, 7, 1, 5$. β) $12, 4, 8, 6$. γ) $10, 3, 3, 7, 3, 9$.

\ex Σωστό ή λάθος; α) Η μέση τιμή είναι πάντα μία από τις τιμές των δεδομένων. β) Αν τα δεδομένα είναι έξι, η διάμεσος είναι ο μέσος όρος της τρίτης και της τέταρτης τιμής, μετά τη σειρά. γ) Για να βρεις τη διάμεσο, πρώτα βάζεις τα δεδομένα σε αύξουσα σειρά. δ) Αν όλες οι τιμές είναι ακέραιοι, η μέση τιμή είναι κι αυτή ακέραιος. ε) Η μέση τιμή μπορεί να είναι μεγαλύτερη από τη μεγαλύτερη τιμή.

\ex Συμπλήρωσε τον πίνακα και βρες τη μέση τιμή: \par
\begin{center}
\renewcommand{\arraystretch}{1.25}
\begin{tabular}{lcccc}
\toprule
Τιμή & $0$ & $1$ & $2$ & $3$\\
Συχνότητα & $4$ & $6$ & $8$ & $2$\\
Τιμή $\times$ συχνότητα & $\ldots$ & $\ldots$ & $\ldots$ & $\ldots$\\
\bottomrule
\end{tabular}
\end{center}

\zoni{Το χρησιμοποιώ}{σπορ, χωριό, σχολείο}

\ex \textbf{Οι βολές.} Ο Νίκος σε έξι προπονήσεις έβαλε $6$, $7$, $5$, $8$, $9$ και $7$ βολές από τις $10$. Βρες τη μέση τιμή και τη διάμεσο.

\ex \textbf{Τα τριαντάφυλλα.} Η θεία Μαρία έχει $9$ τριανταφυλλιές με μέσο όρο $12$ τριαντάφυλλα η καθεμία. α) Πόσα τριαντάφυλλα έχουν όλες μαζί; β) Μια δέκατη τριανταφυλλιά έχει $22$ τριαντάφυλλα. Ποιος είναι ο νέος μέσος όρος;

\ex \textbf{Ο βαθμός που λείπει.} Η Ζωή έγραψε $16$, $18$, $15$, $17$ στα τέσσερα πρώτα τεστ. Τι βαθμό χρειάζεται στο πέμπτο ώστε ο μέσος όρος των πέντε να είναι $17$;

\ex \textbf{Τα βιβλία του καλοκαιριού.} Η Αγγελική ρώτησε $24$ παιδιά πόσα βιβλία διάβασαν το καλοκαίρι.\par
\begin{center}
\renewcommand{\arraystretch}{1.25}
\begin{tabular}{lcccccc}
\toprule
Βιβλία & $0$ & $1$ & $2$ & $3$ & $4$ & $5$\\
Συχνότητα & $1$ & $4$ & $7$ & $7$ & $4$ & $1$\\
\bottomrule
\end{tabular}
\end{center}
Βρες τη μέση τιμή.

\ex \textbf{Το χαρτζιλίκι.} Ο Στέλιος ρώτησε $20$ συμμαθητές πόσα ευρώ παίρνουν την εβδομάδα.\par
\begin{center}
\renewcommand{\arraystretch}{1.25}
\begin{tabular}{lcccc}
\toprule
Ευρώ & $0$--$5$ & $5$--$10$ & $10$--$15$ & $15$--$20$\\
Συχνότητα & $3$ & $8$ & $6$ & $3$\\
\bottomrule
\end{tabular}
\end{center}
Βρες τα κέντρα των κλάσεων και τη μέση τιμή.

\zoni{Το σκέφτομαι}{συγκρίσεις και αιτιολογήσεις}

\ex Πέντε δεδομένα έχουν μέση τιμή $10$. Προσθέτουμε ένα έκτο δεδομένο ίσο με $10$. Πώς αλλάζει η μέση τιμή; Αν το έκτο ήταν $16$, ποια θα ήταν;

\ex Βρες πέντε αριθμούς που να έχουν διάμεσο $7$ και μέση τιμή $10$. Είναι μοναδική η απάντηση;

\ex Το τμήμα Α έχει $10$ μαθητές με μέσο όρο $14$ και το τμήμα Β έχει $20$ μαθητές με μέσο όρο $17$. Ο Στέλιος λέει: «Ο μέσος όρος και των δύο μαζί είναι $\frac{14+17}{2}=15{,}5$». Έχει δίκιο; Βρες τον σωστό μέσο όρο.

\ex Σε ένα μαγαζί του χωριού πέντε εργαζόμενοι παίρνουν $900$, $1000$, $1100$, $1000$ και $5000$ ευρώ τον μήνα (ο τελευταίος είναι ο ιδιοκτήτης). Βρες τη μέση τιμή και τη διάμεσο. Ποια από τις δύο περιγράφει καλύτερα τον «τυπικό» μισθό και γιατί;

\zoni{Για το σπίτι}{μετρήσεις στο σπίτι}

\ex Κράτησε σημειώσεις για τις ώρες που κοιμήθηκες επτά βραδιές. Βρες τη μέση τιμή και τη διάμεσο.

\ex Ρώτησε $10$ ανθρώπους πόσα κινητά τηλέφωνα έχουν στο σπίτι τους. Φτιάξε πίνακα συχνοτήτων και βρες τη μέση τιμή.
